\MathBookSourcePage{339}
\subsection{混合变分形式的例子}
\label{sec:ch12-examples-mixed-formulations}
\setcounter{thmcount}{0}
\setcounter{equation}{0}

高度黏性流体稳态流动的 Stokes 方程为
\MBSourceItem{1}
\begin{equation}
\MathBookFormulaID{formula-ch12-stokes-equations}
\label{eq:ch12-stokes-equations}
-\Delta\chvec u+\chvec{\operatorname{grad}}p=\chvec f,
\qquad
\operatorname{div}\chvec u=0
\quad\text{in }\Omega\subset\mathbb R^n,
\end{equation}
其中 $\chvec u$ 表示流体速度, $p$ 表示压力. 这里使用上一章引入的下划线记号表示向量、向量算子与矩阵. 但在某些情形下也使用推广到三维的显然约定. 为了在所有情形下清楚地表示维数, 使用 $H^s(\Omega)^n$ 代替 $\chvec H^s(\Omega)$, 使用 $L^2(\Omega)^n$ 代替 $\chvec L^2(\Omega)$, 等等.

\MathBookSourcePage{340}
方程~\eqref{eq:ch12-stokes-equations} 表示随后讨论的 Navier--Stokes 方程~\MBUnresolvedReference{eq:ch13-source-13-4-1}{13.4.1} 在 Reynolds 数为零时的极限情形. 对 $\chvec f$ 取非零通常并不常见, 但方程的线性允许把更常见的非齐次边界数据化为齐次数据. 因而, 例如, 假设在边界上具有 Dirichlet 边界条件 $\chvec u=0$.

分部积分可对适当的 $\chvec v$ 和 $q$ 导出变分恒等式
\MBSourceItem{2}
\begin{equation}
\MathBookFormulaID{formula-ch12-stokes-first-variational-identity}
\label{eq:ch12-stokes-first-variational-identity}
a(\chvec u,\chvec v)+b(\chvec v,p)
=\int_\Omega\chvec f\cdot\chvec v\,dx,
\end{equation}
\MBSourceItem{3}
\begin{equation}
\MathBookFormulaID{formula-ch12-stokes-second-variational-identity}
\label{eq:ch12-stokes-second-variational-identity}
b(\chvec u,q)=0,
\end{equation}
其中
\MBSourceItem{4}
\begin{equation}
\MathBookFormulaID{formula-ch12-stokes-a-form}
\label{eq:ch12-stokes-a-form}
a(\chvec u,\chvec v):=
\int_\Omega\sum_{i=1}^n
\chvec{\operatorname{grad}}u_i\cdot
\chvec{\operatorname{grad}}v_i\,dx,
\end{equation}
\MBSourceItem{5}
\begin{equation}
\MathBookFormulaID{formula-ch12-divergence-b-form}
\label{eq:ch12-divergence-b-form}
b(\chvec v,q):=-\int_\Omega(\operatorname{div}\chvec v)q\,dx.
\end{equation}
令
\[
\MathBookFormulaID{formula-ch12-stokes-spaces}
V:=\mathring H^1(\Omega)^n,
\qquad
\Pi:=\left\{q\in L^2(\Omega):\int_\Omega q\,dx=0\right\}.
\]
众所周知(Girault 与 Raviart 1979, Temam 1984), 下列问题有唯一解: 给定 $F\in V'$, 求 $\chvec u\in V$ 和 $p\in\Pi$, 使得
\MBSourceItem{6}
\begin{equation}
\MathBookFormulaID{formula-ch12-stokes-abstract-problem}
\label{eq:ch12-stokes-abstract-problem}
\begin{aligned}
a(\chvec u,\chvec v)+b(\chvec v,p)&=F(\chvec v)&&\forall\chvec v\in V,\\
b(\chvec u,q)&=0&&\forall q\in\Pi.
\end{aligned}
\end{equation}
该问题的适定性部分来自 $a(\cdot,\cdot)$ 在 $\mathring H^1(\Omega)^n$ 上的强制性(参见习题~\ref{exr:ch12-coercivity-stokes-a}).

占据区域 $\Omega$ 的多孔介质中的流体流动模型具有~\MBUnresolvedReference{eq:ch05-source-5-6-6}{5.6.6} 的形式, 即
\MBSourceItem{7}
\begin{equation}
\MathBookFormulaID{formula-ch12-porous-medium-pressure-equation}
\label{eq:ch12-porous-medium-pressure-equation}
-\sum_{i,j=1}^n\frac{\partial}{\partial x_i}
\left(a_{ij}(x)\frac{\partial p}{\partial x_j}(x)\right)
=f(x)\quad\text{in }\Omega,
\end{equation}
其中 $p$ 是压力. 为简单起见, 这里仍取非齐次右端项. Darcy 定律假定流体速度 $\chvec u$ 与 $p$ 的梯度之间满足
\[
\MathBookFormulaID{formula-ch12-darcy-law}
\sum_{j=1}^n a_{ij}(x)\frac{\partial p}{\partial x_j}(x)
=u_i(x),\qquad i=1,\ldots,n.
\]
假设系数 $a_{ij}$ 几乎处处组成对称正定系统, 并与介质的孔隙率有关. 当然, 许多其他物理模型也具有~\eqref{eq:ch12-porous-medium-pressure-equation} 的形式.

\MathBookSourcePage{341}
令 $A(x)$ 表示系数矩阵 $(a_{ij})$ 的几乎处处有定义的逆矩阵, 并写 $\chvec{\operatorname{grad}}p=A\chvec u$. 可对~\eqref{eq:ch12-porous-medium-pressure-equation} 导出形如~\eqref{eq:ch12-stokes-abstract-problem} 的变分形式. 定义
\MBSourceItem{8}
\begin{equation}
\MathBookFormulaID{formula-ch12-porous-medium-a-form}
\label{eq:ch12-porous-medium-a-form}
a(\chvec u,\chvec v):=
\int_\Omega\sum_{i,j=1}^n A_{ij}u_i v_j\,dx,
\end{equation}
而 $b(\cdot,\cdot)$ 仍由~\eqref{eq:ch12-divergence-b-form} 给出. 则~\eqref{eq:ch12-porous-medium-pressure-equation} 的解满足
\MBSourceItem{9}
\begin{equation}
\MathBookFormulaID{formula-ch12-porous-medium-mixed-problem}
\label{eq:ch12-porous-medium-mixed-problem}
\begin{aligned}
a(\chvec u,\chvec v)+b(\chvec v,p)&=0&&\forall\chvec v\in V,\\
b(\chvec u,q)&=F(q)&&\forall q\in\Pi,
\end{aligned}
\end{equation}
其中 $\Pi=L^2(\Omega)$ 且
\[
\MathBookFormulaID{formula-ch12-hdiv-space}
V:=\{\chvec v\in L^2(\Omega)^n:\operatorname{div}\chvec v\in L^2(\Omega)\}.
\]
后一空间称为 $H(\operatorname{div})$, 其自然范数为
\MBSourceItem{10}
\begin{equation}
\MathBookFormulaID{formula-ch12-hdiv-norm}
\label{eq:ch12-hdiv-norm}
\|\chvec v\|_{H(\operatorname{div})}^2
=\|\chvec v\|_{L^2(\Omega)^n}^2
+\|\operatorname{div}\chvec v\|_{L^2(\Omega)}^2
\end{equation}
(参见习题~\ref{exr:ch12-hdiv-hilbert}). 与前一问题不同, 形式 $a(\cdot,\cdot)$ 在整个 $V$ 上并不强制, 因为它只控制 $V$ 的 $L^2$ 范数. 但它在散度为零的函数组成的关键子空间上强制. 在下面的抽象设置中, 这一子空间的作用将会变得清楚.
